\documentclass[11pt,a4paper,parskip=half]{scrartcl}

\usepackage[T1]{fontenc}
\usepackage[utf8]{inputenc}
\usepackage{lmodern}
\usepackage{microtype}
\usepackage[a4paper,left=2.2cm,right=2.2cm,top=2.4cm,bottom=2.6cm,headsep=10pt]{geometry}
\usepackage{xcolor}
\usepackage[table]{xcolor}
\usepackage{amsmath}
\usepackage{amssymb}
\usepackage{booktabs}
\usepackage{array}
\usepackage{tabularx}
\usepackage{colortbl}
\usepackage{enumitem}
\usepackage{tikz}
\usetikzlibrary{positioning,calc,shapes.geometric,backgrounds,patterns,arrows.meta}
\usepackage[most]{tcolorbox}
\usepackage{fontawesome5}
\usepackage[headsepline,footsepline]{scrlayer-scrpage}
\usepackage{caption}
\usepackage{lastpage}
\usepackage{hyperref}

% ---------- palette ----------
\definecolor{fwnavy}{HTML}{1B2A4A}
\definecolor{fwteal}{HTML}{0E7C86}
\definecolor{fwamber}{HTML}{C97A1F}
\definecolor{fwgray}{HTML}{F2F4F7}
\definecolor{fwline}{HTML}{D8DEE6}

\addtokomafont{section}{\color{fwnavy}\sffamily}
\addtokomafont{subsection}{\color{fwnavy}\sffamily}
\addtokomafont{subsubsection}{\color{fwteal}\sffamily}
\addtokomafont{disposition}{\color{fwnavy}}

\hypersetup{
  colorlinks=true,
  linkcolor=fwteal,
  urlcolor=fwteal,
  citecolor=fwteal,
  pdftitle={SY32L-S1 Firmware Technical Reference},
  pdfauthor={Synapse Microsystems -- Firmware Platform Team}
}

% ---------- page style ----------
\clearpairofpagestyles
\ihead{\small\sffamily\color{fwnavy} SY32L-S1 Firmware Technical Reference}
\ohead{\small\sffamily\color{fwteal} FPT-2026-0114}
\ifoot{\small\sffamily\color{gray} Synapse Microsystems -- Confidential}
\ofoot{\small\sffamily\color{gray} Page \pagemark\ of \pageref{LastPage}}
\setkomafont{headsepline}{\color{fwline}}
\setkomafont{footsepline}{\color{fwline}}

\setlength{\parindent}{0pt}

% ---------- small helpers ----------
\newtcolorbox{noticebox}[1][]{%
  enhanced,
  colback=fwgray,
  colframe=fwteal,
  boxrule=0.6pt,
  left=8pt,right=8pt,top=5pt,bottom=5pt,
  arc=1.5pt,
  fonttitle=\sffamily\bfseries,
  coltitle=fwnavy,
  #1
}

\newtcolorbox{warnbox}[1][]{%
  enhanced,
  colback=white,
  colframe=fwamber,
  boxrule=0.7pt,
  left=8pt,right=8pt,top=5pt,bottom=5pt,
  arc=1.5pt,
  fonttitle=\sffamily\bfseries,
  coltitle=fwamber,
  title={\faExclamationTriangle\ Firmware Note},
  #1
}

\newcolumntype{Y}{>{\raggedright\arraybackslash}X}

\begin{document}

% ============================================================
% TITLE BLOCK
% ============================================================
\begin{tikzpicture}[remember picture,overlay]
  \fill[fwnavy] ($(current page.north west)$) rectangle ($(current page.north west)+(\paperwidth,-3.1cm)$);
  \fill[fwteal] ($(current page.north west)+(0,-3.1cm)$) rectangle ($(current page.north west)+(\paperwidth,-3.16cm)$);
\end{tikzpicture}

\vspace*{-1.55cm}
\begin{minipage}{0.68\textwidth}
  \color{white}\sffamily
  {\fontsize{11}{13}\selectfont \faFile\ TECHNICAL REFERENCE MANUAL}\\[6pt]
  {\fontsize{22}{26}\selectfont\bfseries SY32L-S1 Microcontroller}\\[2pt]
  {\fontsize{15}{18}\selectfont Firmware Register \& Interrupt Reference}
\end{minipage}
\hfill
\begin{minipage}{0.27\textwidth}
  \raggedleft\color{white}\sffamily\small
  Synapse Microsystems\\
  Firmware Platform Team\\[4pt]
  \faLock\ Confidential -- NDA Distribution
\end{minipage}

\vspace{0.9cm}

\begin{center}
\begin{tabularx}{\textwidth}{|Y|Y|Y|Y|}
\hline
\rowcolor{fwgray}
\textbf{\small Document ID} & \textbf{\small Revision} & \textbf{\small Applies To} & \textbf{\small Status} \\
\hline
\small FPT-2026-0114 & \small v2.3.0 & \small SY32L-S1, Rev.\ C silicon & \small Released \\
\hline
\end{tabularx}
\end{center}

\vspace{0.4cm}

% ============================================================
% SCOPE
% ============================================================
\section*{Scope and Audience}
\addcontentsline{toc}{section}{Scope and Audience}

This document is the firmware-facing register and interrupt reference for the
Synapse SY32L-S1 microcontroller family. It is written for embedded engineers
who are bringing up board support packages, writing peripheral drivers, or
debugging interrupt latency on Rev.\ C silicon. It does not cover electrical
characteristics, package mechanicals, or qualification data -- those live in
the SY32L-S1 Datasheet (FPT-2026-0110). Register bit layouts, reset values,
and access types below are normative for firmware and take precedence over
any inline code comments in the reference driver package.

Readers are assumed to be comfortable with ARM Cortex-M0+ programming,
memory-mapped I/O, and nested vectored interrupt controllers. Address values
are given in hexadecimal without separators (e.g.\ \texttt{0x08000000}).
Bit ranges are written \texttt{[MSB:LSB]}; a single-bit field is written with
one index only.

\begin{noticebox}[title={\faInfoCircle\ Register Notation}]
\small
\textbf{RW} read/write \quad
\textbf{RO} read-only \quad
\textbf{RC} read-clears (write ignored) \quad
\textbf{W1C} write-1-to-clear \quad
\textbf{Res.} reserved, must be written as reset value.
Undefined behaviour results from writing reserved bits to anything other
than their reset value.
\end{noticebox}

% ============================================================
% ARCHITECTURE
% ============================================================
\section{Device Architecture}

The SY32L-S1 pairs a Cortex-M0+ core running at up to 64\,MHz with a
single-cycle AHB matrix and two APB peripheral bridges. Firmware talks to
every peripheral in this document through memory-mapped registers on APB1
or APB2; there are no port-mapped I/O instructions on this core.

\begin{center}
\begin{tikzpicture}[
  node distance=6mm and 10mm,
  block/.style={draw=fwnavy, thick, rounded corners=1.5pt, fill=fwgray,
    minimum height=9mm, minimum width=28mm, align=center, font=\small\sffamily},
  core/.style={block, fill=fwnavy, text=white, minimum width=36mm},
  bus/.style={draw=fwteal, thick, -{Latex[length=2mm]}},
]
  \node[core] (cpu) {Cortex-M0+ Core\\[-1pt]\footnotesize 64\,MHz, NVIC};
  \node[block, above left=6mm and -2mm of cpu] (flash) {Flash\\[-1pt]\footnotesize 128\,KB};
  \node[block, above right=6mm and -2mm of cpu] (sram) {SRAM\\[-1pt]\footnotesize 32\,KB};

  \node[block, below=14mm of cpu, minimum width=110mm, fill=white] (bridge) {AHB / APB1 / APB2 Bridge Matrix};

  \node[block, below left=10mm and -26mm of bridge, minimum width=24mm] (gpio) {GPIO};
  \node[block, right=4mm of gpio, minimum width=24mm] (usart) {USART1/2};
  \node[block, right=4mm of usart, minimum width=24mm] (tim) {TIM1--TIM3};
  \node[block, right=4mm of tim, minimum width=24mm] (adc) {ADC1};
  \node[block, right=4mm of adc, minimum width=24mm] (rcc) {Clock Ctrl};

  \draw[bus] (flash.south) -- (cpu.north -| flash.south);
  \draw[bus] (sram.south) -- (cpu.north -| sram.south);
  \draw[bus] (cpu.south) -- (bridge.north);
  \draw[bus] (bridge.south -| gpio.north) -- (gpio.north);
  \draw[bus] (bridge.south -| usart.north) -- (usart.north);
  \draw[bus] (bridge.south -| tim.north) -- (tim.north);
  \draw[bus] (bridge.south -| adc.north) -- (adc.north);
  \draw[bus] (bridge.south -| rcc.north) -- (rcc.north);
\end{tikzpicture}
\end{center}

\subsection{Memory Map}

\begin{center}
\begin{tabularx}{\textwidth}{@{}l l l Y@{}}
\toprule
\textbf{Region} & \textbf{Start} & \textbf{End} & \textbf{Notes} \\
\midrule
Boot ROM (System) & 0x1FFF0000 & 0x1FFF3FFF & Factory bootloader, read-only, entered when \texttt{BOOT0}=1 \\
Option Bytes & 0x1FFFF800 & 0x1FFFF80F & Read protection, watchdog-on-reset, BOR level \\
Main Flash & 0x08000000 & 0x0801FFFF & 128\,KB application code, erase page = 1\,KB \\
SRAM & 0x20000000 & 0x20007FFF & 32\,KB, includes stack and .bss \\
APB1 Peripherals & 0x40000000 & 0x40007FFF & TIM2--TIM3, USART2, I2C1 \\
APB2 Peripherals & 0x40010000 & 0x40017FFF & GPIO, USART1, TIM1, ADC1, EXTI \\
AHB Peripherals & 0x40020000 & 0x40023FFF & DMA1, RCC (clock control), CRC \\
Cortex-M0+ Private Bus & 0xE000E000 & 0xE000EFFF & NVIC, SysTick, SCB \\
\bottomrule
\end{tabularx}
\end{center}

\newpage

% ============================================================
% CLOCK CONTROL REGISTER
% ============================================================
\section{Peripheral Registers}

Base addresses for each peripheral instance are given in
Table~\ref{tab:regmap}; all offsets below are relative to that base.

\begin{center}
\begin{tabularx}{\textwidth}{@{}l l Y@{}}
\toprule
\textbf{Peripheral} & \textbf{Base Address} & \textbf{Instance Notes} \\
\midrule
RCC (Clock Control) & 0x40021000 & Single instance, AHB \\
GPIOA & 0x40010800 & 16 pins, APB2 \\
GPIOB & 0x40010C00 & 16 pins, APB2 \\
USART1 & 0x40013800 & APB2, up to 4\,Mbit/s \\
USART2 & 0x40004400 & APB1, up to 2\,Mbit/s \\
TIM1 & 0x40012C00 & Advanced, APB2 \\
TIM2, TIM3 & 0x40000000, 0x40000400 & General-purpose, APB1 \\
ADC1 & 0x40012400 & 12-bit, APB2 \\
\bottomrule
\end{tabularx}
\captionof{table}{Peripheral base addresses.}
\label{tab:regmap}
\end{center}

\subsection{RCC\_CFGR -- Clock Control Register (offset 0x04)}

Reset value: \texttt{0x00000000}. Controls the PLL multiplier, AHB
prescaler, and reports which oscillator currently drives the system clock.

\begin{center}
\small
\begin{tabular}{|*{16}{c|}}
\hline
15 & 14 & 13 & 12 & 11 & 10 & 9 & 8 & 7 & 6 & 5 & 4 & 3 & 2 & 1 & 0 \\
\hline
\multicolumn{4}{c|}{PLLMUL} & \multicolumn{3}{c|}{AHBPRE} & HSEON & HSERDY & HSION & HSIRDY & \multicolumn{3}{c|}{SWS} & \multicolumn{2}{c|}{SW} \\
\hline
\end{tabular}
\end{center}

\begin{center}
\begin{tabularx}{\textwidth}{@{}l l l l Y@{}}
\toprule
\textbf{Field} & \textbf{Bits} & \textbf{Access} & \textbf{Reset} & \textbf{Description} \\
\midrule
PLLMUL & [15:12] & RW & 0x0 & PLL multiplier, output = HSI/2 $\times$ (PLLMUL+2), clamped to 64\,MHz \\
AHBPRE & [11:9] & RW & 0x0 & AHB prescaler: 000=/1 ... 111=/512 \\
HSEON & [8] & RW & 0 & Enable external 8\,MHz crystal oscillator \\
HSERDY & [7] & RO & 0 & Set by hardware once HSE is stable \\
HSION & [6] & RW & 1 & Enable internal 8\,MHz RC oscillator \\
HSIRDY & [5] & RO & 1 & Set by hardware once HSI is stable \\
SWS & [4:2] & RO & 0x0 & System clock actually in use: 000=HSI, 001=HSE, 010=PLL \\
SW & [1:0] & RW & 0x0 & Requested system clock switch, mirrored into SWS once locked \\
\bottomrule
\end{tabularx}
\end{center}

\begin{warnbox}
Firmware must poll \texttt{HSERDY} (or \texttt{HSIRDY}) before writing
\texttt{SW}. Switching to an oscillator that has not reported ready leaves
\texttt{SWS} at its previous value and silently ignores the write -- this
has been the root cause of two field returns where the bootloader appeared
to hang for exactly the HSE startup timeout.
\end{warnbox}

\subsection{USART\_CR1 -- Control Register 1 (offset 0x0C)}

Reset value: \texttt{0x00000000}. Applies to both USART1 and USART2.

\begin{center}
\small
\begin{tabular}{|*{16}{c|}}
\hline
15 & 14 & 13 & 12 & 11 & 10 & 9 & 8 & 7 & 6 & 5 & 4 & 3 & 2 & 1 & 0 \\
\hline
OVER8 & \multicolumn{2}{c|}{Res.} & M & WAKE & PCE & PS & PEIE & TXEIE & TCIE & RXNEIE & IDLEIE & TE & RE & RWU & SBK \\
\hline
\end{tabular}
\end{center}

\begin{center}
\begin{tabularx}{\textwidth}{@{}l l l l Y@{}}
\toprule
\textbf{Field} & \textbf{Bits} & \textbf{Access} & \textbf{Reset} & \textbf{Description} \\
\midrule
OVER8 & [15] & RW & 0 & 0: 16$\times$ oversampling, 1: 8$\times$ oversampling (higher baud rates) \\
M & [12] & RW & 0 & Word length: 0 = 8 data bits, 1 = 9 data bits \\
PCE & [10] & RW & 0 & Parity control enable \\
PS & [9] & RW & 0 & Parity selection: 0 = even, 1 = odd \\
TXEIE & [7] & RW & 0 & Interrupt when TDR is empty \\
TCIE & [6] & RW & 0 & Interrupt on transmission complete \\
RXNEIE & [5] & RW & 0 & Interrupt when RDR has unread data \\
TE & [3] & RW & 0 & Transmitter enable \\
RE & [2] & RW & 0 & Receiver enable \\
SBK & [0] & RW & 0 & Send break; cleared by hardware once the break has been sent \\
\bottomrule
\end{tabularx}
\end{center}

\subsection{TIMx\_CR1 -- Timer Control Register 1 (offset 0x00)}

Reset value: \texttt{0x00000000}. Applies to TIM1, TIM2, and TIM3.

\begin{center}
\small
\begin{tabular}{|*{16}{c|}}
\hline
15 & 14 & 13 & 12 & 11 & 10 & 9 & 8 & 7 & 6 & 5 & 4 & 3 & 2 & 1 & 0 \\
\hline
\multicolumn{6}{c|}{Res.} & \multicolumn{2}{c|}{CKD} & ARPE & \multicolumn{2}{c|}{CMS} & DIR & OPM & URS & UDIS & CEN \\
\hline
\end{tabular}
\end{center}

\begin{center}
\begin{tabularx}{\textwidth}{@{}l l l l Y@{}}
\toprule
\textbf{Field} & \textbf{Bits} & \textbf{Access} & \textbf{Reset} & \textbf{Description} \\
\midrule
CKD & [9:8] & RW & 0x0 & Clock division for dead-time and filter sampling \\
ARPE & [7] & RW & 0 & Auto-reload preload enable; buffers ARR writes until next update \\
CMS & [6:5] & RW & 0x0 & Center-aligned mode selection (TIM1 only; 0x0 on TIM2/3) \\
DIR & [4] & RW & 0 & Counter direction: 0 = up, 1 = down \\
OPM & [3] & RW & 0 & One-pulse mode: counter stops at next update event \\
URS & [2] & RW & 0 & Update request source restricted to overflow/underflow \\
UDIS & [1] & RW & 0 & Update event disable \\
CEN & [0] & RW & 0 & Counter enable \\
\bottomrule
\end{tabularx}
\end{center}

\newpage

% ============================================================
% INTERRUPT VECTOR TABLE
% ============================================================
\section{Interrupt Vector Table}

The vector table below is placed at \texttt{0x08000000} at reset (or
relocated via \texttt{SCB\_VTOR} for the bootloader's application slot at
\texttt{0x08002000}). Priorities are 2-bit (0 = highest) as implemented by
this core's reduced NVIC.

\begin{center}
\small
\begin{tabularx}{\textwidth}{@{}c l l c Y@{}}
\toprule
\textbf{\#} & \textbf{Offset} & \textbf{Exception / IRQ} & \textbf{Prio} & \textbf{Description} \\
\midrule
--- & 0x00 & Initial Stack Pointer & --- & Loaded into \texttt{MSP} before Reset\_Handler runs \\
1 & 0x04 & Reset & $-3$ & Fixed highest priority, not configurable \\
2 & 0x08 & NMI & $-2$ & Clock security system failure routes here \\
3 & 0x0C & HardFault & $-1$ & Bus/usage faults with no dedicated handler \\
11 & 0x2C & SVCall & configurable & Used by the reference RTOS for context switch \\
14 & 0x38 & PendSV & configurable & RTOS tail-chaining hook \\
15 & 0x3C & SysTick & configurable & 1\,ms tick, drives the software timebase \\
16 & 0x40 & WWDG & 1 & Window watchdog early-warning \\
17 & 0x44 & PVD & 2 & Programmable voltage detector threshold crossed \\
18 & 0x48 & RTC & 1 & Alarm or wakeup from the always-on domain \\
19 & 0x4C & FLASH & 2 & End-of-write / end-of-erase from the flash controller \\
20 & 0x50 & RCC & 2 & Oscillator or PLL ready \\
21 & 0x54 & EXTI0\_1 & 0 & External lines 0 and 1 \\
22 & 0x58 & EXTI2\_3 & 0 & External lines 2 and 3 \\
23 & 0x5C & EXTI4\_15 & 0 & External lines 4 through 15 \\
25 & 0x64 & DMA1\_CH1 & 1 & DMA channel 1 complete/half/error \\
28 & 0x70 & ADC1 & 1 & End-of-conversion, watchdog, overrun \\
29 & 0x74 & TIM1\_UP & 0 & TIM1 update event \\
31 & 0x7C & TIM2 & 0 & TIM2 global interrupt \\
32 & 0x80 & TIM3 & 0 & TIM3 global interrupt \\
34 & 0x88 & I2C1 & 1 & Combined event and error \\
35 & 0x8C & SPI1 & 1 & TX/RX buffer and error \\
36 & 0x90 & USART1 & 0 & TX empty, RX not empty, framing/overrun errors \\
37 & 0x94 & USART2 & 0 & Same as USART1, second instance \\
\bottomrule
\end{tabularx}
\end{center}

\begin{noticebox}[title={\faInfoCircle\ Shared Handlers}]
\small
\texttt{EXTI4\_15} covers twelve GPIO lines. Drivers must read
\texttt{EXTI\_PR} inside the handler and clear only the bit they own --
clearing the whole register has caused missed edges on adjacent pins in the
motor-control reference design.
\end{noticebox}

% ============================================================
% BOOT SEQUENCE
% ============================================================
\section{Boot Sequence}

\begin{minipage}{0.58\textwidth}
\begin{enumerate}[label=\textbf{\arabic*.}, leftmargin=1.6em, itemsep=3pt]
  \item Power-on reset holds the core until \texttt{VDD} crosses the
    brown-out threshold latched in the option bytes.
  \item Boot ROM samples \texttt{BOOT0}. High selects the system
    bootloader at \texttt{0x1FFF0000}; low proceeds to the application.
  \item Core fetches the initial stack pointer from offset
    \texttt{0x00} of the active vector table, then the Reset\_Handler
    address from offset \texttt{0x04}.
  \item \texttt{SystemInit()} configures HSE, starts the PLL, and only
    then switches \texttt{RCC\_CFGR.SW} -- see the notice in
    Section~2.1.
  \item \texttt{.data} is copied from flash to SRAM and \texttt{.bss}
    is zeroed by the startup file.
  \item \texttt{main()} is called with interrupts still globally
    masked; application code unmasks them once peripheral clocks and
    NVIC priorities are configured.
  \item Independent watchdog, if enabled by option bytes, must receive
    its first refresh within 26\,ms of this point.
\end{enumerate}
\end{minipage}
\hfill
\begin{minipage}{0.38\textwidth}
\begin{tcolorbox}[colback=fwgray,colframe=fwnavy,boxrule=0.5pt,arc=1.5pt,
  title={Cold Boot Timing},fonttitle=\sffamily\bfseries\small,
  coltitle=white,colbacktitle=fwnavy]
\small
\begin{tabular}{@{}l r@{}}
Reset to HSI stable & 2\,\textmu s \\
HSE crystal startup & 1.8\,ms \\
PLL lock & 90\,\textmu s \\
Flash wait-state calc & 4\,\textmu s \\
\texttt{main()} entry & 2.1\,ms \\
\end{tabular}
\end{tcolorbox}
\end{minipage}
\par

\newpage

% ============================================================
% POWER MODES
% ============================================================
\section{Power Modes}

\begin{center}
\begin{tabularx}{\textwidth}{@{}l l l Y@{}}
\toprule
\textbf{Mode} & \textbf{Typ. Current} & \textbf{Wake Latency} & \textbf{Wake Sources} \\
\midrule
Run & 4.8\,mA @ 64\,MHz & --- & --- \\
Sleep & 1.6\,mA & 1 cycle & Any enabled interrupt \\
Stop & 3.2\,\textmu A & 4\,\textmu s & EXTI lines, RTC alarm, USART wake \\
Standby & 0.4\,\textmu A & 60\,\textmu s & WKUP pins, RTC alarm, IWDG reset only \\
\bottomrule
\end{tabularx}
\end{center}

\begin{warnbox}
SRAM contents are lost in Standby. Any driver state that must survive
(calibration constants, communication sequence counters) has to be pushed
to the 128-byte backup register block before entering this mode -- there
is no automatic retention.
\end{warnbox}

Entry to Stop or Standby is via \texttt{PWR\_CR} followed by the Cortex-M0+
\texttt{WFI} instruction. Firmware should disable the SysTick interrupt
before \texttt{WFI} in Stop mode; a pending SysTick otherwise wakes the
core every 1\,ms and defeats the purpose of the low-power entry.

% ============================================================
% REVISION HISTORY
% ============================================================
\section{Revision History}

\begin{center}
\begin{tabularx}{\textwidth}{@{}l l Y@{}}
\toprule
\textbf{Rev.} & \textbf{Date} & \textbf{Change} \\
\midrule
v1.0.0 & 2025-02-10 & Initial release, Rev.\ A silicon \\
v2.0.0 & 2025-09-04 & Added TIM1 advanced-timer fields, Rev.\ B silicon \\
v2.2.0 & 2026-03-18 & Corrected \texttt{RCC\_CFGR.SWS} reset behaviour after HSE fallback \\
v2.3.0 & 2026-07-01 & Rev.\ C silicon: added \texttt{EXTI4\_15} shared-handler notice, updated Stop-mode current \\
\bottomrule
\end{tabularx}
\end{center}

\vspace{0.6cm}
\begin{center}
\small\color{gray}
Synapse Microsystems -- Firmware Platform Team \hspace{4pt}|\hspace{4pt} FPT-2026-0114 \hspace{4pt}|\hspace{4pt} Confidential, NDA distribution only
\end{center}

\end{document}
